\Askhsh{29784}{2}
\begin{ylh}{1}
    \item 9.4. Γενίκευση του Πυθαγόρειου θεωρήματος
    \item 10.1. Πολυγωνικά χωρία
    \item 10.4. Άλλοι τύποι για το εμβαδόν τριγώνου.
\end{ylh}
\begin{wrapfigure}[18]{r}{7cm}
    \centering
    \begin{tikzpicture}[scale=0.8]
        \tkzDefPoint(0,0){A}
        \tkzDefPoint(5,0){B}
        % Define D using polar coordinates relative to A for angle DAB = 120 degrees and AD = 3
        \tkzDefPointWith[polar=3:120](A){D}
        
        % To define G such that angle DBG = 120 degrees and BG = 8:
        % Rotate point D around B by -120 degrees (clockwise) to find the direction for BG.
        \tkzDefPointBy[rotation=center B angle -120](D) \tkzGetPoint{P_rotated}
        % G is on the line BP_rotated at distance 8 from B.
        \tkzDefPointWith[euclidian,K=8](B,P_rotated){G}
        
        % Draw the quadrilateral and the diagonal BD
        \tkzDrawSegments(A,B B,G G,D D,A B,D)
        
        % Label points
        \tkzLabelPoint[below](A){A}
        \tkzLabelPoint[below](B){B}
        \tkzLabelPoint[right](G){$\Gamma$}
        \tkzLabelPoint[above left](D){$\Delta$}
        
        % Label side lengths
        \tkzLabelSegment[below](A,B){5}
        \tkzLabelSegment[left](A,D){3}
        \tkzLabelSegment[right](B,G){8}
        
        % Mark and label angles
        \tkzMarkAngle[arc,size=0.8cm,fill=gray!30](B,A,D)
        \tkzLabelAngle[pos=1.1](B,A,D){$120^\circ$}
        
        \tkzMarkAngle[arc,size=0.8cm,fill=gray!30](D,B,G)
        \tkzLabelAngle[pos=1.1](D,B,G){$120^\circ$}
    \end{tikzpicture}
\end{wrapfigure}

Στο τετράπλευρο του παρακάτω σχήματος έχουμε για τις πλευρές $AB=5$, $B\Gamma=8$, $A\Delta=3$ και για τις γωνίες $\Delta\widehat{A}B = 120^\circ$ όπως και $\Delta\widehat{B}\Gamma = 120^\circ$.
\begin{alist}
    \item Να δείξετε ότι το μήκος της $B\Delta$ είναι 7 μονάδες. \monades{8}
    \item Να υπολογίσετε το μήκος της $\Gamma\Delta$. \monades{8}
    \item Να υπολογίσετε το εμβαδόν του τριγώνου $\Gamma B \Delta$. \monades{9}
\end{alist}